


\subsection{Rule-based and Programmatic Reasoning}
\label{subsec:p3_s7_1_5}
% - Discrete task logic and sequential decision-making: using rules/operators to specify when and in what order actions should be performed.
% - FSM, PDDL, STRIPS, Behavior Trees, If-then rules
% - Applications: multi-step task planning, logic-guided assembly, agent behavior switching

Rule-based and programmatic reasoning is the explicit, human-auditable layer that decides \emph{when} an action is allowed, \emph{in what order} actions happen, and \emph{what to do next} when conditions change, using symbolic operators, transition rules, or program-like control structures. In a robot it is the supervisory layer that sequences reusable skills (navigate, grasp, insert) and invokes geometric solvers or learned policies only when their preconditions hold, which is what makes a decision explainable (why the detour when the corridor was blocked), verifiable, and able to carry commonsense and safety constraints through a long task. Figure~\ref{fig:p3_s7_1_5_rule_programmatic_abstract} gives the abstract formulation.

\begin{figure}[t]
    \centering
    \includegraphics[width=0.99\linewidth]{figures/part3/section7/715_program.pdf}
    \caption{Abstract task formulation of rule-based and programmatic reasoning: (i) symbolic task structure, (ii) four representations for executing decision logic, and (iii) executable plans, reactive policies, and monitored execution traces with replanning after a disturbance.}
    \label{fig:p3_s7_1_5_rule_programmatic_abstract}
\end{figure}

\subsubsection{Task Definition}
A symbolic planning problem is a tuple $\Pi=\langle \mathcal{S},\mathcal{A},s_0,G\rangle$ with symbolic states $\mathcal{S}$, an operator library $\mathcal{A}$, an initial state $s_0$, and a goal $G$; each operator $a\in\mathcal{A}$ carries a precondition and effects,
\begin{equation}
a=\big(\mathrm{pre}(a),\ \mathrm{eff}^{+}(a),\ \mathrm{eff}^{-}(a)\big),\qquad s'=\big(s\setminus\mathrm{eff}^{-}(a)\big)\cup\mathrm{eff}^{+}(a),
\end{equation}
and is applicable in $s$ iff $\mathrm{pre}(a)\subseteq s$. The planning output is a sequence $\pi=\langle a_1,\dots,a_k\rangle$ whose successive application to $s_0$ reaches a state satisfying $G$; the reactive variant instead returns a policy $\mu$ that maps the current observation and symbolic state to a behavior. Depending on the formalism, the inputs also carry resources, temporal constraints, and exogenous events.

\subsubsection{Methodology}
The formalisms split by whether they \emph{plan} (search for a correct sequence before acting) or \emph{react} (choose a behavior continually at run time); deployed systems usually do both, with a planner proposing the sequence and an execution structure (FSM, behavior tree, rule engine) monitoring conditions and triggering recovery.

\paragraph{Classical Symbolic Planning (STRIPS/PDDL and Heuristic Search).}
STRIPS~\cite{Fikes1971STRIPS} fixed the representation: planning as search over symbolic states with action schemas, so that a plan can be checked against the model and the continuous motion layer sits behind a clean abstraction boundary. PDDL~\cite{McDermott1998PDDL} standardized domain and problem specification so that planners could be compared across navigation, logistics, and assembly, with extensions (ADL, numeric fluents, temporal constructs) that capture robotic requirements without abandoning explicit semantics. Graphplan~\cite{Blum1997PlanningGraph} builds a planning graph whose mutual-exclusion constraints prune incompatible action combinations early, and FF~\cite{Hoffmann2001FF} made forward search fast with relaxed-plan heuristics that ignore delete effects, which capture ordering constraints such as unstacking before stacking in object-rich domains. Task-and-motion planning (§\ref{subsec:p3_s7_1_4}) extends these planners with geometric feasibility checks~\cite{Kaelbling2011HierarchicalNow,Garrett2021TAMPReview,Zhao2024TAMPsurvey}; the common pattern is to keep the symbolic reasoning explicit, delegate feasibility to motion and IK solvers, and wrap the plan in an execution structure with monitors and recovery.

\paragraph{Finite-State Machines and Reactive Mode Switching.}
An FSM switches among a few well-defined behaviors (navigate, align, grasp) with explicit transition guards, which makes it transparent and fast but prone to state explosion as contingencies multiply. Statecharts~\cite{Harel1987Statecharts} add hierarchy and concurrency, so that complex controllers are nested reusable submachines with precise semantics; Brooks' subsumption architecture~\cite{Brooks1986Subsumption} layers behaviors so that higher ones override lower ones, a blueprint for composing navigation and avoidance without global planning, systematized as behavior-based robotics~\cite{Arkin1998BehaviorBased}; teleo-reactive programs~\cite{Nilsson1994TeleoReactive} evaluate conditions continuously and fire the highest-priority satisfied rule; and ROS tooling such as FlexBE~\cite{conner2023flexbe} provides hierarchical, event-driven state machines at deployment scale.

\paragraph{Behavior Trees.}
A behavior tree re-evaluates its conditions on every tick and chooses among structured alternatives (sequence, fallback, parallel), so that long-horizon execution survives disturbances without committing to one fixed sequence; subtrees encapsulate reusable routines, guards and fallbacks trigger immediate recovery, and the composition semantics fit perception and motion skills. BTs originated in game AI~\cite{Isla2005Halo2}; Marzinotto \etal~\cite{Marzinotto2014BT} formalized their semantics and their relation to other switching structures, and Colledanchise and {\"O}gren~\cite{Colledanchise2017BT,Colledanchise2018BTbook} showed which compositions preserve safety and robustness and consolidated the design principles. Compiling or wrapping a symbolic plan into a BT keeps the nominal task structure while exposing runtime guards and recovery branches~\cite{mart2021behaviourtree}; ARIAC systems pair PDDL planning with BT execution monitors to handle failures, occlusions, and defective parts~\cite{liu2024autonomous}; and event-driven BTs~\cite{Agis2020EDBT} react to non-periodic stimuli rather than waiting for the next tick.

\paragraph{Logic-based Execution and Production Rule Systems.}
A production system keeps a working memory of facts and fires forward-chaining rules when their conditions hold, which suits event-condition-action policies, mode switching, and safety enforcement (no-go zones, mutual exclusions, priorities) and depends on fast incremental matching. OPS5~\cite{Forgy1981OPS5} set the model; the Rete algorithm~\cite{Forgy1982Rete} caches partial matches so that large rule sets update in milliseconds as facts change; CLIPS~\cite{Riley1991CLIPS} made the toolchain portable; and Soar~\cite{Laird1987Soar} integrated large rule bases with planning and execution in one architecture. In robotics, KnowRob~\cite{tenorth2013knowrob} and CRAM~\cite{beetz2010cram,beetz2023cram} join symbolic knowledge to rule-driven execution for reasoning about actions, effects, and everyday manipulation, and LTLMoP~\cite{finucane2010ltlmop} compiles temporal-logic specifications into correct-by-construction reactive controllers whose safety and sequencing constraints are explicit at the program level.

\subsubsection{Benchmarks}
The International Planning Competition supplies the canonical PDDL domains and problems across classical, temporal and numeric, probabilistic, and HTN tracks, compared by plan quality, makespan, and runtime~\cite{McDermott1998PDDL}, with \texttt{planning.domains}~\cite{PlanningDomains} mirroring and extending the sets for reproducible evaluation. Long-horizon embodied suites test execution under partial observability and variation: BEHAVIOR-100 in iGibson and BEHAVIOR-1K in OmniGibson~\cite{Srivastava2022BEHAVIOR,li2024behavior1k} specify household activities at the predicate level, ALFRED~\cite{shridhar2020alfred} and TEACh~\cite{Padmakumar2022TEACh} score instruction-conditioned action sequences, and RoboTHOR~\cite{Deitke2020RoboTHOR} and the Habitat challenge~\cite{szot2021habitat2} pair high-level decision logic with navigation policies under SPL. ARIAC~\cite{liu2024autonomous} injects faulty parts, occlusions, and failures into kitting and assembly, and FurnitureBench~\cite{Heo2023FurnitureBench} and RoboCup@Home~\cite{Wisspeintner2009RoboCupHome} stress recovery and explainability on real tasks; qualitative solvers over RCC-8~\cite{Randell1992RCC} and Allen relations~\cite{janhunen2019allen} are tested on consistency problems, and composite pipelines reuse the TAMP and OMPL harnesses of §\ref{subsec:p3_s7_1_4}.

\subsubsection{Challenges and Trends}
The results on these benchmarks divide along the closed-world boundary. On IPC-style domains with fully specified preconditions and effects, classical planners remain strong; on long-horizon embodied tasks the limit is perception noise and grounding ambiguity rather than search, and a logically applicable action may be geometrically unrealizable because of a collision, an IK failure, or a missing grasp. Hand-authored operators, rules, and recovery logic also grow combinatorially with object count and contingency. Integrated task-and-motion planning (§\ref{subsec:p3_s7_1_4}) answers the coupling problem by interleaving symbolic search with feasibility checks or compiling symbolic commitments into optimization back-ends~\cite{LozanoPerez2014ConstraintTAMP,garrett2020pddlstream,toussaint2015lgp}, and hierarchical abstraction and compositional skill models contain the search while keeping it auditable~\cite{Garrett2021TAMPReview,Zhao2024TAMPsurvey}.

Robustness to noise and nonstationarity is why the industrial and instruction-conditioned benchmarks reward compositions rather than single formalisms: a behavior tree around a symbolic plan makes disturbances and recoveries explicit through runtime guards, fallbacks, and safety-preserving composition~\cite{Marzinotto2014BT,Colledanchise2017BT,Colledanchise2018BTbook}, and Rete-style engines~\cite{Forgy1982Rete} supply fast event-condition-action policies over streaming facts, the pattern ARIAC entries converge on~\cite{liu2024autonomous}.

Learned perception and language models increasingly propose the high-level procedure or heuristic while the symbolic structure stays as the safety shell that enforces preconditions, effects, and constraints: SayCan~\cite{ahn2022icanisay} grounds language in robot skills by value-based scoring, and plan-to-BT compilation~\cite{mart2021behaviourtree} exposes parallelism, guards, and recovery at run time. Domain modeling, heuristic guidance, and recovery synthesis are being automated piece by piece while the programmatic structure is retained for auditing, debugging, and verification.
